\documentclass[11pt]{article}
\usepackage{ghb}

% On the site these become semantic block quotations styled by site.css.  The
% PDF gets a matching shaded panel rather than a plain indented quotation.
\ifdefined\ghbHTML
  \newenvironment{promptbox}{\begin{quote}}{\end{quote}}
\else
  \usepackage{tcolorbox}
  \newtcolorbox{promptbox}{
    colback=black!4,
    colframe=black!20,
    boxrule=0.4pt,
    arc=2pt,
    left=8pt,
    right=8pt,
    top=6pt,
    bottom=6pt,
    before skip=10pt,
    after skip=10pt
  }
\fi

\title{I guess we live in the $P = NP$ world after all}
\date{2026-09-18}

\summary{AI haunts Gödel's Lost Letter}

\begin{document}

\begin{abstract}
One of the most important questions in computer science is whether it is intrinsically
more time-consuming to \textit{come up} with an answer than to \textit{verify} one.
This is the so-called $P = NP$ question: if they're equal it's as ``easy'' to come up with an answer as it is to verify the answer.
That would be pretty weird: doesn't coming up with answers require more creativity? 
So most computer scientists are pretty sure $P \neq NP$. But maybe not! It's a \textsc{big} question.
Just recently it seemed to me, \textit{maybe it doesn't matter anymore?}
\end{abstract}

\section{A Magic Answer Box}

Imagine we had a machine that could answer any question. Better than that though, it gives you an 
incontrovertible \textit{proof} backing its answer up.

A machine like that would be easy to build: It's just a computer! It doesn't even need any fancy AI.
Here's how it works: You give it a (mathematical) statement $F$, and some ginormous number, $n$. 
It tells you whether there is a proof of $F$ of length at most $n$.
If there is, it can hand you the proof.

And it's easy! The magic answer box just generates all possible strings of length at most $n$, and checks if any of them are valid proofs of $F$. 
Voil\`a! Creativity is automated.

\section{Complexity}

The only issue with the magic answer box is that it's \textit{slow}. The only general algorithm we know that is guaranteed to work is very, very slow. 
Sure, \textit{checking} one is easier. But we don't know an efficient algorithm for finding the proofs. It all sounds like the $P = NP$ question! 

As one famous meat-mathematician once wrote to another: 

\begin{promptbox}
\large\itshape
``The mental work of a mathematician concerning Yes-or-No questions
could be completely replaced by a machine.''

\par\medskip
	\hfill\normalsize --- \href{https://rjlipton.com/the-gdel-letter}{Kurt G\"odel's ``lost'' letter to John von Neumann, 1956}
\end{promptbox}

So, $P = NP$ being true would be crazy - we might then have an efficient algorithm for automating mathematics! 

\section{One Million Dollars}

With the recent smashing of Navier-Stokes by \href{https://openai.com/index/navier-stokes-solution/}{OpenAI}, I've had the
\href{https://www.claymath.org/millennium-problems/}{Millennium Prize Problems} on my mind.
They serve as dramatic lighthouses for mathematical endeavour.
These are the amazing, rock hard problems that you get a million dollars for settling (with a proof!). 

The difficulty with having them on my mind though, is that I don't know what any of them really \textit{mean}. Except,
of course for the $P = NP$ question. That one, I know a little about. In recent weeks, I've heard people
say that it is likely the \textit{hardest} of them all. Even the magician mathematician Terence 
Tao is \href{https://www.youtube.com/watch?v=PtsrAw1LR3E&t=2727s}{on record} saying that it would probably be the last problem to fall.

So I've dared to wonder recently: maybe I'll know for sure whether $P \neq NP$ soon.

\section{Brainworm}

Of course, you can guess the dreadful imposition that comes blacking in one's mind: We \textit{have} automated creativity.
An AI has settled Navier-Stokes. I personally have had AIs answer many publishable questions in computer science. 
I can't get it out of my head, it turns out, \textit{we live in the $P = NP$ world}!

\section{Of Course Not Quite}

Now, sure, yes, $P \neq NP$. Or maybe not. And G\"odel's magic-answer-box machine is \textit{exhaustive}, and AI is not.
Maybe that makes all the difference. The AIs might fail on some problems, and the magic-answer-box would be guaranteed not to ``miss'' any proof. 
But it does make me wonder. Perhaps every natural mathematical question whose answer is within reach of human understanding is also within reach of AI.
Maybe \textit{of course}, and not perhaps at all. 

\section{I See a Darkness}

So why do I care if $P = NP$ anymore? It feels like the only millennium problem I understood, is maybe just another 
curiosity now. Perhaps the funniest world would
be the one where AIs automate creativity, and quickly prove $P = NP$, via an $O(n^{100})$ time algorithm for \href{https://en.wikipedia.org/wiki/Satisfiability}{SAT}.

\end{document}
